\subsection{Transformer Models}\label{s:method_rq1}

\subsubsection{Task}

We train a modified auto-regressive transformer model to map from $[d_1, d_2, \texttt{SEP}, \texttt{MASK}, \texttt{MASK}]$ to $[d_1, d_2, \texttt{SEP}, o_1, o_2]$, where $d_1$ and $d_2$ are uniformly sampled symbols from $[0, 99]$ and $(o_1, o_2)$ are the outputs in the same order as the input: $(o_1, o_2) = (d_1, d_2)$. We refer to these as symbols rather than integers to emphasise that numeric value is irrelevant; integers are used to represent an arbitrary set of distinct symbols.
\texttt{MASK} and \texttt{SEP} are always represented by out-of-distribution symbols $100, 101$ respectively, hence the model should not confuse these with input symbols.

\begin{table}[h]
\centering
\caption{
Our custom attention mask for input $[d_1, d_2, \texttt{SEP}, o_1, o_2]$, where \xmark~means the attention weight was set to $-\infty$ before softmax, and \cmark~means it was computed normally. 
The $d_i$ tokens attend only to other $d_i$ and \texttt{SEP}, while the $o_i$ tokens attend causally to the separator and previous output tokens.
This prevents direct copying of input embeddings from $d_i$ into $o_i$, forcing the model to learn a compressed representation at \texttt{SEP}.
Self-attention is enabled only for $d_1$ to prevent numerical instability during training and does not affect task performance.
}
\label{tab:mask}
\vspace{0.5em} % This command adds a small vertical space
\begin{tabular}{@{}cccccc@{}}
\toprule
& $d_1$ & $d_2$ & \texttt{SEP} & $o_1$ & $o_2$ \\
\midrule
$d_1$ & \cmark & \xmark & \xmark & \xmark & \xmark \\
$d_2$ & \cmark & \xmark & \xmark & \xmark & \xmark \\
\texttt{SEP} & \cmark & \cmark & \xmark & \xmark & \xmark \\
$o_1$ & \xmark & \xmark & \cmark & \xmark & \xmark \\
$o_2$ & \xmark & \xmark & \cmark & \cmark & \xmark \\
\bottomrule
\end{tabular}
\end{table}

The attention pattern is constrained so that output positions cannot attend to and simply copy from the input positions: $o_1$ can only attend to \texttt{SEP} and $o_2$ can only attend to $o_1$ and \texttt{SEP} (Table \ref{tab:mask}). This forces the model to use a compressed representation of the input bigram at \texttt{SEP} to complete the task.
Loss and accuracy is computed on only $o_1$ and $o_2$, and the final outputs for $d_1$, $d_2$ and \texttt{SEP} are ignored. In this way, 50\% accuracy can be achieved by correctly predicting all $o_1$ but failing all $o_2$ (or vice versa).

Intuitively, the model should therefore learn a mechanism such that, in the first layer, it forms a compressed representation of two fixed-length vectors $(d_1, d_2)$ and stores it at \texttt{SEP} (single fixed-length vector), and in the second layer, it recovers the original $(d_1, d_2)$ from that compressed representation and stores them in $(o_1, o_2)$, preserving input order.
In this way, we structure the task to force the model to learn a mechanism that requires relational composition \cite{wattenberg2024relational}.

We chose this narrow setting to isolate the problem solved by the subpath pre-computation in \citep{brinkmann2024mechanistic}, and simplified to an ordered bigram to find a minimal list copying circuit. 

\subsubsection{Model specification and training}

Our training dataset consists of 80\% of all possible inputs (8000 of the total 100*100 $(d_1, d_2)$ bigrams). The test set is the remaining 20\% of unseen inputs. We optimise the mean cross-entropy loss over the output positions:
\begin{equation*}
  \mathcal{L} = -\frac{1}{N}\sum_{i=1}^{N}
    \log P_\theta(o_i \mid \texttt{SEP},\, o_1, \ldots, o_{i-1})
\end{equation*}
where $P_\theta$ is the model distribution with parameters $\theta$ and $N$ is the $N$-gram size (number of $d_i$ in the input, which equals the number of $o_i$ in the output). In our case $N=2$ (a bigram).

\begin{figure}[htbp]
  \centering
  \includegraphics[width=\linewidth]{figures/listlen-nlayer-heatmap.pdf}
  \caption{
  Results from our grid search across $2 \leq n_\text{layers} \leq 8$ and $N$-gram size $2 \leq N \leq 8$.
  Each point represents the average accuracy across 30 differently seeded training runs with 50,000 steps each. Other parameters were fixed: $d_\text{model}=64$, $n_\text{heads}=1$, and no LayerNorm, bias, MLP, or value or output matrices. We hypothesise that the drop-off in accuracy for larger $n_\text{layers}$ results from the insufficiency of 50,000 steps to train that many parameters.
}
  \label{fig:listlen_nlayer_heatmap}
\end{figure}

We train several transformers with different numbers of layers, $n_\text{layers} \in [2,8]$, on this task (Figure \ref{fig:listlen_nlayer_heatmap}). 
Models with $n_\text{layers} \geq 3$ often achieve 100\% accuracy for bigrams ($N=2$), as the additional layer(s) allow the model to directly copy between token positions, circumventing the composition task, as shown for a three-layer model in Figure \ref{fig:attn-flow}. 

Whilst the three-layer model is more accurate, 
we focus on two-layer transformers, which provide strong incentive to learn relational composition at the \texttt{SEP} token after layer 1, as alternatives (direct copying) require additional layers. We select a two-layer model with test accuracy 91.45\%, which lies between the mean and maximum across 30 seeds ($0.8607 \pm 0.1080$, max $0.9405$). We prefer a near-mean run over the single best seed so that our claims can better characterise what the model class typically learns, not an outlier best solution (Table \ref{tab:params-grid-search}).

This reflects previous work. \citet{brinkmann2024mechanistic} find that models with $L$ layers trained to solve pathfinding problems in $N$-depth binary trees create the path one step per layer if $L > N$, and only use the subpath compression mechanism if $L-1 < N$. 
In our model, we meet the same constraint where $N$ corresponds to the input $N$-gram, rather than tree depth. For a bigram $(d_1, d_2)$, we thus require a two-layer model.
However, we additionally investigate the three-layer model shown in Figure \ref{fig:attn-flow} since it is the smallest model that achieves 100\% accuracy, which provides a useful mechanistic comparison.

\begin{table}
\caption{
Grid search results where $n_\text{layers}=2$ and $n_\text{heads}=1$ are fixed. F = False, T = True. 
We write LN for LayerNorm, and $W_V$ and $W_O$ are the value and output matrices respectively, where False means the matrix was frozen to the identity during training. 
We further test the effect of adding an MLP to each layer since our models are attention-only.
We additionally ablate the residual stream dimension $d_\text{model}$ and $n_\text{heads}$ in the second and third sections respectively.
Accuracy reports the mean $\pm$ standard deviation of test accuracy across 30 seeds, with Max being the best single run. Each run was 50,000 training steps. The best result in each section is \textbf{bold} and the selected configuration used in this paper is shown in the first row for easier comparison.
}
\label{tab:params-grid-search}
\centering
\setlength{\tabcolsep}{3pt}
\begin{tabular}{c c c c c c c l c}
\toprule
$d_\text{model}$ & $n_\text{heads}$ & LN & Bias & $W_V$ & $W_O$ & MLP & Mean Accuracy & Max \\
\midrule
64 & 1 & F & F & F & F & F & 0.8607 $\pm$ 0.1080 & 0.9405 \\
\midrule
64 & 1 & F & F & F & T & F & 0.8921 $\pm$ 0.0545 & 0.9210 \\
64 & 1 & F & F & T & F & F & 0.8668 $\pm$ 0.1079 & 0.9205 \\
64 & 1 & F & F & T & T & F & 0.7438 $\pm$ 0.1801 & 0.8960 \\
64 & 1 & F & T & F & F & F & 0.8652 $\pm$ 0.1144 & 0.9370 \\
64 & 1 & F & T & F & T & F & 0.8455 $\pm$ 0.1388 & 0.9125 \\
64 & 1 & F & T & T & F & F & 0.8559 $\pm$ 0.1248 & 0.9210 \\
64 & 1 & F & T & T & T & F & 0.8012 $\pm$ 0.1619 & 0.8948 \\
64 & 1 & T & F & F & F & F & 0.9308 $\pm$ 0.0045 & 0.9395 \\
64 & 1 & T & F & F & T & F & 0.9212 $\pm$ 0.0048 & 0.9340 \\
64 & 1 & T & F & T & F & F & 0.9203 $\pm$ 0.0036 & 0.9290 \\
64 & 1 & T & F & T & T & F & 0.9242 $\pm$ 0.0054 & 0.9335 \\
64 & 1 & T & T & F & F & F & 0.9241 $\pm$ 0.0052 & 0.9323 \\
64 & 1 & T & T & F & T & F & 0.9212 $\pm$ 0.0048 & 0.9310 \\
64 & 1 & T & T & T & F & F & 0.9210 $\pm$ 0.0040 & 0.9305 \\
64 & 1 & T & T & T & T & F & 0.9248 $\pm$ 0.0050 & 0.9343 \\
64 & 1 & F & F & F & F & T & 0.8966 $\pm$ 0.0102 & 0.9180 \\
64 & 1 & F & F & F & T & T & 0.8893 $\pm$ 0.0102 & 0.9080 \\
64 & 1 & F & F & T & F & T & 0.8888 $\pm$ 0.0085 & 0.9145 \\
64 & 1 & F & F & T & T & T & 0.8315 $\pm$ 0.0221 & 0.8590 \\
64 & 1 & F & T & F & F & T & 0.8978 $\pm$ 0.0125 & 0.9293 \\
64 & 1 & F & T & F & T & T & 0.8870 $\pm$ 0.0096 & 0.9018 \\
64 & 1 & F & T & T & F & T & 0.8901 $\pm$ 0.0094 & 0.9110 \\
64 & 1 & F & T & T & T & T & 0.8335 $\pm$ 0.0200 & 0.8638 \\
64 & 1 & T & F & F & F & T & 0.9408 $\pm$ 0.0079 & 0.9595 \\
64 & 1 & T & F & F & T & T & 0.9485 $\pm$ 0.0068 & 0.9597 \\
64 & 1 & T & F & T & F & T & \textbf{0.9498 $\pm$ 0.0053} & 0.9597 \\
64 & 1 & T & F & T & T & T & 0.9497 $\pm$ 0.0083 & 0.9630 \\
64 & 1 & T & T & F & F & T & 0.9415 $\pm$ 0.0073 & \textbf{0.9653} \\
64 & 1 & T & T & F & T & T & 0.9487 $\pm$ 0.0062 & 0.9610 \\
64 & 1 & T & T & T & F & T & 0.9489 $\pm$ 0.0072 & 0.9595 \\
64 & 1 & T & T & T & T & T & 0.9497 $\pm$ 0.0086 & 0.9613 \\
\midrule
256 & 1 & F & F & F & F & F & 0.9225 $\pm$ 0.0098 & 0.9340 \\
128 & 1 & F & F & F & F & F & \textbf{0.9304 $\pm$ 0.0044} & \textbf{0.9375} \\
32 & 1 & F & F & F & F & F & 0.5626 $\pm$ 0.1345 & 0.9052 \\
8 & 1 & F & F & F & F & F & 0.3521 $\pm$ 0.1183 & 0.6340 \\
\midrule
64 & 2 & F & F & F & F & F & \textbf{0.8256 $\pm$ 0.1399} & \textbf{0.9565} \\
64 & 4 & F & F & F & F & F & 0.4822 $\pm$ 0.0416 & 0.5795 \\
64 & 8 & F & F & F & F & F & 0.3788 $\pm$ 0.0806 & 0.5045 \\
\bottomrule
\end{tabular}
\end{table}

After a comprehensive hyperparameter sweep for the two-layer model (Table \ref{tab:params-grid-search}), we find that we can ablate attention bias and LayerNorm \cite{ba2016layer} and freeze the attention value and output matrices to the identity while maintaining near top performance (achieved: 0.940 accuracy, top: 0.965 accuracy). This reduces attention outputs to weighted sums of the previous residual stream, which allows us to more easily analyse outputs and learned mechanisms that use attention.
We set $d_{model} = 64$, which is the residual stream dimension, which achieved the best top accuracy out of the tested $d_{model}$ values.
We use single-headed attention which has the highest mean test accuracy out of tested $n_\text{heads}$ combinations and it's much simpler to analyse than multi-headed attention \cite{elhage2021mathematical}.

\begin{table}[htbp]
\caption{Final model and training configuration.}
\label{tab:model_config}
\centering
\begin{tabular}{lp{3.8cm}}
\toprule
\textbf{Parameter} & \textbf{Value} \\
\midrule
Number of layers & 2 or 3 \\
Number of heads & 1 \\
Residual stream dimension ($d_{\text{model}}$) & 64 \\
Attention head dimension ($d_{\text{head}}$) & 64 \\
LayerNorm & None \\
Bias & None \\
Value matrix ($W_V$) & Identity (frozen) \\
Output matrix ($W_O$) & Identity (frozen) \\
MLP & None \\
\midrule
Learning rate & $1 \times 10^{-3}$ \\
Optimiser & \texttt{AdamW} \\
Betas & $(0.9,\ 0.999)$ \\
Weight decay & $0.01$ \\
Batch size & 2048 \\
LR schedule & None \\
Max training steps & 50{,}000 \\
% Early stopping & val accuracy $\geq 99.9\%$ \\
Mixed precision & FP16 (AMP) \\
\midrule
Hardware & NVIDIA GeForce RTX 2080 Ti \\
\bottomrule
\end{tabular}
\end{table}

Table \ref{tab:model_config} shows our final configurations.
Overall, we have 29,760 trainable parameters for the two-layer model, and 125,056 for the three-layer model.
These are the token and positional embedding matrices, the unembedding matrix, and the key and query matrices from the attention layers. 

Whilst these are substantial constraints, it is normal for toy models in the mechanistic interpretability literature to use similar constraints to provide a controlled experimental environment for analysis \citep{nanda2023progress, elhage2021mathematical}.

\subsubsection{Evaluation Methodology}

To validate that the model is achieving the task as intended, we ablate each attention pattern probability and measure impact on performance. This highlights which parts of the model's attention are causally significant for performance.
To successfully complete the task with strong performance, we predict that attention from \texttt{SEP} to the input bigram $(d_1, d_2)$ will be critical in the first layer, and the attention from output bigram $(o_1, o_2)$ to \texttt{SEP} will be critical in the second layer.

To interpret the algorithm learned by the model to complete the task, we perform circuits analysis \cite{ameisen2025circuit} by deriving formulae for the logits at $(o_1, o_2)$.
This demonstrates specifically how the model uses and composes different information for computation, and this analysis depth is only feasible on a constrained toy model as described.
